\section{Introduction}
\label{sec:introduction}

Low Earth orbit (LEO) satellites can complement terrestrial networks
to provide Internet-of-Things (IoT) services in remote areas.
To support access from a large number of IoT devices, grant-free random
access (GFRA) allows active devices to transmit without first obtaining
a grant, thereby reducing signaling overhead and access latency.
However, the receiver has no prior knowledge of which devices are
active. Meanwhile, limited pilot resources make it impractical to assign
mutually orthogonal pilots to all devices. The receiver must therefore
perform joint activity detection and channel estimation (JADCE) from
nonorthogonal pilots \cite{liu2018}. In LEO satellite links, this task
is further complicated by large propagation delays and Doppler shifts.

To compensate for these effects before uplink transmission, the new
radio non-terrestrial network (NR-NTN) framework in Release~17 of the
3rd Generation Partnership Project (3GPP) relies on devices equipped
with a global navigation satellite system (GNSS) receiver. Devices use
their GNSS positions, satellite ephemeris, and the common timing advance
broadcast by the network for delay and Doppler pre-compensation
\cite{3gpp38300}. GNSS hardware and frequent pre-compensation increase
the cost and energy consumption of IoT devices. To support devices
without GNSS receivers or access to GNSS services, Release~19 includes
a study of GNSS-independent operation \cite{3gpp22865}. Handling large
differential delays and Doppler shifts at the satellite receiver reduces
the need for user-specific pre-compensation. The differential delay can
then exceed one symbol duration, while the differential Doppler shift
can exceed one subcarrier spacing. Terrestrial orthogonal frequency
division multiplexing (OFDM) GFRA schemes designed for small timing and
frequency offsets and block-fading channels must therefore be adapted
to these conditions.

Orthogonal time-frequency space (OTFS) modulation supports communication
over doubly selective channels through a delay--Doppler representation
that exploits diversity in both time and frequency \cite{hadani2017}.
For access schemes that handle large offsets at the satellite receiver,
the waveform must also permit simple terminal modulation. Affine
frequency division multiplexing (AFDM) uses a one-dimensional discrete
affine Fourier transform (DAFT) with adjustable chirp parameters.
Properly chosen parameters separate multipath components in the DAFT
domain and enable their diversity to be exploited. AFDM modulation
requires an inverse fast Fourier transform and chirp multiplications,
providing a suitable implementation for terminals with limited
computational resources \cite{bemani2022}.

The channel representation is central to JADCE with AFDM. Conventional
GFRA receivers exploit sparse user activity through compressed sensing
and use channel structure to improve recovery. For example, the
massive multiple-input multiple-output (MIMO) receiver in
\cite{ke2020} jointly exploits activity sparsity and virtual
angular-domain channel sparsity. In GNSS-independent satellite access,
Shen et al.\ \cite{shen2025} combine multi-frame OTFS transmission with
message passing (MP) to jointly recover device activity, channel coefficients, and data
symbols, exploiting channel sparsity in the delay--Doppler--angle domain.
For LEO AFDM GFRA, the graph neural
network in \cite{ding2026} refines a sparse channel prior using nodes
associated with delay--Doppler grid points. Physical delays and Doppler
shifts, however, are continuous. Fractional offsets spread power across
grid coefficients and weaken the assumed sparsity. Increasing the grid
resolution enlarges the dictionary, while estimating continuous offsets
adds path parameters coupled with the unknown gains. Direct estimation
of the AFDM effective channel offers an alternative that incorporates
these offsets into the channel to be recovered.

Score-based generative models learn channel distributions from samples
and use this information for MIMO channel estimation \cite{arvinte2023}.
For AFDM channel estimation, \cite{huang2026} learns the score of the
effective channel and performs posterior inference using the AFDM
observation model. This representation captures fractional-delay and
fractional-Doppler effects without prescribing sparse, grid-aligned
channel support. In multiuser access, the learned channel information
must also support activity detection and interference suppression.
The turbo MP receiver in \cite{cai2026} integrates
score-based channel denoising with linear estimation for MIMO-OFDM
JADCE. It uses the learned prior through local denoising updates,
avoiding online reverse diffusion sampling. These results motivate
combining a learned AFDM effective-channel prior with the joint
estimation structure of MP.

Effective use of this prior also depends on the inputs produced during
message passing. Score pretraining uses Gaussian-perturbed channels,
whereas finite system dimensions and structured sensing matrices can
produce denoiser inputs that depart from the equivalent Gaussian model
\cite{cai2026}. Damping smooths successive messages, but strong damping
can require more iterations to reach a given accuracy
\cite{yoshida2023}. In JADCE, channel and activity errors are coupled
through the feedback: the activity probability weights the channel
estimate used in subsequent interference suppression, and the channel
estimate provides evidence for activity detection. Thus the design
problem includes both adapting the denoiser to its received inputs and
coordinating the two estimators across successive updates.

Deep unfolding turns an iterative receiver into a trainable sequence of
estimation modules \cite{he2020}. We use this structure to optimize
JADCE from received AFDM pilots over a prescribed number of MP layers.
Each layer contains a channel denoiser, an activity classifier, and
trainable message parameters. Joint training propagates channel and
activity losses through their shared feedback: activity supervision
adjusts channel denoising, while channel losses at later layers refine
the activity estimates used by earlier layers. This provides a direct
training mechanism for the interaction between channel recovery,
activity decisions, and multiuser interference suppression.

In this paper, we develop a multi-frame AFDM transmission scheme and a
deep-unfolded MP receiver for GNSS-independent MIMO-AFDM satellite IoT
access. The main contributions are summarized as follows.
\begin{enumerate}[label=\arabic*)]
\item We derive a multi-frame AFDM input--output model for large
differential delays and Doppler shifts, retaining both integer and
fractional offsets. Separating the inter-frame Doppler phase evolution
places these phase factors in the pilot operator and allows all pilot
frames within an access block to share an effective-channel matrix.
The resulting formulation combines additional pilot observations
without introducing a separate set of channel coefficients for each
frame.

\item We develop a channel and activity estimation module based on a
learned multiantenna AFDM effective-channel prior. The score-based
denoiser estimates each user's channel conditional on being active, and a
trainable classifier estimates its activity probability from the
denoiser input and channel estimate. Their product is used to form the
feedback to the linear estimator, integrating off-grid channel
structure and soft activity information into multiuser MP.

\item We unfold the MP receiver into layers with independently
trainable denoisers, activity classifiers, and message parameters.
A progressive training procedure adds layers and jointly fine-tunes
the receiver using channel and activity supervision. By propagating
both losses through the inter-layer messages, the training adapts the
two estimators to the evolving receiver inputs and their effects on
subsequent interference suppression.
\end{enumerate}

% Numerical findings will be added after independent evaluation of the new receiver.
Section~\ref{sec:system} presents the system model.
Sections~\ref{sec:algorithm}--\ref{sec:overall} develop the receiver,
its training procedure, and computational complexity. Sections~VI and
VII present numerical results and conclusions, respectively.
